%\title{Plants in Teop material culture}
\title{The material culture of the Teop people of Bougainville}
%\subtitle{Texts on plants, house building, canoes and fishing}
\author{Ulrike Mosel}
\renewcommand{\lsISBNdigital}{978-3-96110-594-6}
\renewcommand{\lsISBNhardcover}{978-3-98554-211-6}
\renewcommand{\lsID}{601}
\BookDOI{10.5281/zenodo.22012183}
\typesetter{Christian Döhler}
\proofreader{Matthew Korte}
\renewcommand{\lsSeries}{otc}
\renewcommand{\lsSeriesNumber}{12}
\illustrator{Neville Vitahi; Photographs: Owen Kasinori and Ulrike Mosel}
\dedication{To the memory of Enoch Horai Magum, Naphtali Maion and Neville Vitahi}
\BackBody{In this text collection, nine Teop authors and one Teop graphic artist document the use of palms, trees, bamboo and other plants in their material culture. 

The book consists of six main parts. Following the introductory chapters on the Teop language and the historical context of the text collection in Part I, Parts II to V contain 32 Teop texts with illustrations and English translations on the topics of plants, house-building, canoe-building and traditional fishing techniques. Part VI provides all the Teop texts and their translations, with an additional linguistic analysis and glossing of each word in interlinearized format.

The chapters in Part I and the brief introductions to the texts in Parts II to V are by the editor, who also produced the English translations in collaboration with the Teop authors, as well as the linguistic analysis and glossing in Part VI.}
